\section{Introduction}
\label{sec:introduction}

As the number of processing elements integrated on a single chip increases,
scalable on-chip communication becomes a first-class design concern.
Network-on-Chip (NoC) architectures have emerged as the standard answer to this
challenge, replacing shared buses and point-to-point links with a structured,
packet-switched interconnect fabric~\cite{dally2001}.

Existing NoC designs span a broad design space. Wormhole switching minimizes
buffering by forwarding flits as soon as the header is decoded, but exposes the
network to head-of-line blocking. Virtual channels (VCs) were introduced to
mitigate this blocking~\cite{dally1992}, at the cost of duplicated buffers that
can represent the dominant area contribution of a router~\cite{mello2005}. Circuit
switching provides fully dedicated, contention-free paths but requires path
reservation before any data can flow.

\textbf{HyNoC} takes a different approach: it combines the path-dedication
property of circuit switching---through a source-routed request/grant handshake at
each hop---with the flit-level, FIFO-backed flow control of wormhole switching.
This hybrid model is particularly well suited to FPGA targets, where block RAMs
used as FIFOs are scarce resources, and to distributed VLIW computing workloads,
where communication patterns are largely known at compile time and can therefore be
routed statically to avoid congestion.

Originally inspired by the Hermes NoC~\cite{moraes2004}, HyNoC departs
significantly from its ancestor by introducing distributed arbitration, relative
source routing with multicast support, per-port asynchronous clock domains, and a
parallel round-robin arbiter that grants requests in a one or two cycles.

The contributions of this paper are:
\begin{itemize}
  \item A precise characterization of HyNoC's hybrid switching model and its
        implications for latency and area;
  \item A discussion of static route management as an alternative to virtual
        channels, supported by a combinatorial analysis of shortest-path counts
        in mesh and higher-dimensional topologies;
  \item An analysis of the scope and limitations of the no-VC design choice;
  \item An open-source RTL implementation targeting both simulation and FPGA
        synthesis, complemented by the Veriparse toolkit~\cite{veriparse} for
        source-to-source elaboration of advanced RTL constructs and faster
        simulation;
  \item Verilator co-simulation results measuring the deterministic, linear
        per-hop latency and demonstrating link-disjoint quadrant isolation under
        a LLaMA~3~8B-scale distributed GeMV workload.
\end{itemize}

The remainder of the paper is organized as follows.
Section~\ref{sec:related} reviews related work.
Section~\ref{sec:architecture} describes the HyNoC architecture.
Section~\ref{sec:switching} details the hybrid switching model and the rationale
for avoiding virtual channels.
Section~\ref{sec:target} discusses the target application domain.
Section~\ref{sec:implementation} presents implementation details.
Section~\ref{sec:evaluation} reports simulation results, and
Section~\ref{sec:conclusion} concludes.
